\section{蒙古元的征服、海域外交与环境危机}

\subsection{西夏中亚战争}

攻夏、灭西辽和征花剌子模。战争与外交在，首先改变的往往不是地图颜色，而是道路能否通行、军粮由谁转运、城镇如何开闭、居民往何处迁徙以及地方首领同哪一方交涉。不同语言材料需校年 因此，征服和边界不应被写成一次战役之后立即完成的事实，而应被追踪为军事行动、环境约束和行政接收不断相互调整的过程。

这里的基本证据来自《元史》、\textit{《元朝秘史》}、波斯文和高丽材料、条格、碑刻与环境考古。不同政权的本纪、编年和战报常同意战事发生，却在军数、斩获、攻守责任和谈判条件上彼此冲突。军队报告天然服务奖惩和合法性，后出的亡国或胜利叙事又容易把复杂选择压缩成必然结果。故本书只把能够交叉确认的时间、地点和大致后果写作稳定节点；具体兵力、伤亡、城内情形和个人动机，均保留为需要多方材料核对的问题。

环境并非战争故事的背景布景。河流、山地、草场、渡口、运河、疫病、旱涝与季节决定补给速度，也影响逃亡、安置和复耕的可能。城址、古河道、道路遗存和区域考古能约束文本中的路线，却不能单独量化一次行军或一场灾害的损失。把地理条件同财政和文书联系起来，才可解释同一军令何以在不同地区产生不同的执行结果。

迁徙者、降附者、俘虏、军户、商人和寺院人员并不是被动的战争数字。他们在文书、墓志、碑刻、地方志或宗教材料中偶尔留下姓名与去向，但保存极不均衡。一个家族、一座城市或一批俘虏不能代表全体人口；同样，史书的沉默也不能被理解为没有流动。叙事应分开国家宣布的接收、军队实际控制和居民长期生活的改变。

外交语言也有自己的限制。册封、贡赐、兄弟之国、盟约、降附等称谓说明谈判的仪式形式，却不能直接说明税收、司法或地方认同。对同一条约，应分别追问文本传本、使节往返、边市、堡寨和后续摩擦；对同一征服，应分别追问都城、乡村、边区和海域的时间差。这样写，才能避免把多政权世界误读为单一中枢的扩张史。

攻夏、灭西辽和征花剌子模 的意义正在于使、信息和人群流动短暂可见。材料不足之处，本书不补造完整场景，而以来源差异保留疑问；这既是对战争受难者和地方行动者的谨慎，也是对账本能够承担范围的界定。

\subsection{金朝征服与多方联盟}

攻金、取中都和灭金。战争与外交在，首先改变的往往不是地图颜色，而是道路能否通行、军粮由谁转运、城镇如何开闭、居民往何处迁徙以及地方首领同哪一方交涉。征服叙事不能覆盖居民流动 因此，征服和边界不应被写成一次战役之后立即完成的事实，而应被追踪为军事行动、环境约束和行政接收不断相互调整的过程。

这里的基本证据来自《元史》、\textit{《元朝秘史》}、波斯文和高丽材料、条格、碑刻与环境考古。不同政权的本纪、编年和战报常同意战事发生，却在军数、斩获、攻守责任和谈判条件上彼此冲突。军队报告天然服务奖惩和合法性，后出的亡国或胜利叙事又容易把复杂选择压缩成必然结果。故本书只把能够交叉确认的时间、地点和大致后果写作稳定节点；具体兵力、伤亡、城内情形和个人动机，均保留为需要多方材料核对的问题。

环境并非战争故事的背景布景。河流、山地、草场、渡口、运河、疫病、旱涝与季节决定补给速度，也影响逃亡、安置和复耕的可能。城址、古河道、道路遗存和区域考古能约束文本中的路线，却不能单独量化一次行军或一场灾害的损失。把地理条件同财政和文书联系起来，才可解释同一军令何以在不同地区产生不同的执行结果。

迁徙者、降附者、俘虏、军户、商人和寺院人员并不是被动的战争数字。他们在文书、墓志、碑刻、地方志或宗教材料中偶尔留下姓名与去向，但保存极不均衡。一个家族、一座城市或一批俘虏不能代表全体人口；同样，史书的沉默也不能被理解为没有流动。叙事应分开国家宣布的接收、军队实际控制和居民长期生活的改变。

外交语言也有自己的限制。册封、贡赐、兄弟之国、盟约、降附等称谓说明谈判的仪式形式，却不能直接说明税收、司法或地方认同。对同一条约，应分别追问文本传本、使节往返、边市、堡寨和后续摩擦；对同一征服，应分别追问都城、乡村、边区和海域的时间差。这样写，才能避免把多政权世界误读为单一中枢的扩张史。

攻金、取中都和灭金 的意义正在于使、信息和人群流动短暂可见。材料不足之处，本书不补造完整场景，而以来源差异保留疑问；这既是对战争受难者和地方行动者的谨慎，也是对账本能够承担范围的界定。

\subsection{忽必烈继承战争}

蒙哥去世、争位和李璮叛乱。战争与外交在，首先改变的往往不是地图颜色，而是道路能否通行、军粮由谁转运、城镇如何开闭、居民往何处迁徙以及地方首领同哪一方交涉。继承冲突影响地方征发 因此，征服和边界不应被写成一次战役之后立即完成的事实，而应被追踪为军事行动、环境约束和行政接收不断相互调整的过程。

这里的基本证据来自《元史》、\textit{《元朝秘史》}、波斯文和高丽材料、条格、碑刻与环境考古。不同政权的本纪、编年和战报常同意战事发生，却在军数、斩获、攻守责任和谈判条件上彼此冲突。军队报告天然服务奖惩和合法性，后出的亡国或胜利叙事又容易把复杂选择压缩成必然结果。故本书只把能够交叉确认的时间、地点和大致后果写作稳定节点；具体兵力、伤亡、城内情形和个人动机，均保留为需要多方材料核对的问题。

环境并非战争故事的背景布景。河流、山地、草场、渡口、运河、疫病、旱涝与季节决定补给速度，也影响逃亡、安置和复耕的可能。城址、古河道、道路遗存和区域考古能约束文本中的路线，却不能单独量化一次行军或一场灾害的损失。把地理条件同财政和文书联系起来，才可解释同一军令何以在不同地区产生不同的执行结果。

迁徙者、降附者、俘虏、军户、商人和寺院人员并不是被动的战争数字。他们在文书、墓志、碑刻、地方志或宗教材料中偶尔留下姓名与去向，但保存极不均衡。一个家族、一座城市或一批俘虏不能代表全体人口；同样，史书的沉默也不能被理解为没有流动。叙事应分开国家宣布的接收、军队实际控制和居民长期生活的改变。

外交语言也有自己的限制。册封、贡赐、兄弟之国、盟约、降附等称谓说明谈判的仪式形式，却不能直接说明税收、司法或地方认同。对同一条约，应分别追问文本传本、使节往返、边市、堡寨和后续摩擦；对同一征服，应分别追问都城、乡村、边区和海域的时间差。这样写，才能避免把多政权世界误读为单一中枢的扩张史。

蒙哥去世、争位和李璮叛乱 的意义正在于使、信息和人群流动短暂可见。材料不足之处，本书不补造完整场景，而以来源差异保留疑问；这既是对战争受难者和地方行动者的谨慎，也是对账本能够承担范围的界定。

\subsection{江南接收和对日远征}

南下、崖山与海上远征。战争与外交在，首先改变的往往不是地图颜色，而是道路能否通行、军粮由谁转运、城镇如何开闭、居民往何处迁徙以及地方首领同哪一方交涉。海战损失有战报夸张 因此，征服和边界不应被写成一次战役之后立即完成的事实，而应被追踪为军事行动、环境约束和行政接收不断相互调整的过程。

这里的基本证据来自《元史》、\textit{《元朝秘史》}、波斯文和高丽材料、条格、碑刻与环境考古。不同政权的本纪、编年和战报常同意战事发生，却在军数、斩获、攻守责任和谈判条件上彼此冲突。军队报告天然服务奖惩和合法性，后出的亡国或胜利叙事又容易把复杂选择压缩成必然结果。故本书只把能够交叉确认的时间、地点和大致后果写作稳定节点；具体兵力、伤亡、城内情形和个人动机，均保留为需要多方材料核对的问题。

环境并非战争故事的背景布景。河流、山地、草场、渡口、运河、疫病、旱涝与季节决定补给速度，也影响逃亡、安置和复耕的可能。城址、古河道、道路遗存和区域考古能约束文本中的路线，却不能单独量化一次行军或一场灾害的损失。把地理条件同财政和文书联系起来，才可解释同一军令何以在不同地区产生不同的执行结果。

迁徙者、降附者、俘虏、军户、商人和寺院人员并不是被动的战争数字。他们在文书、墓志、碑刻、地方志或宗教材料中偶尔留下姓名与去向，但保存极不均衡。一个家族、一座城市或一批俘虏不能代表全体人口；同样，史书的沉默也不能被理解为没有流动。叙事应分开国家宣布的接收、军队实际控制和居民长期生活的改变。

外交语言也有自己的限制。册封、贡赐、兄弟之国、盟约、降附等称谓说明谈判的仪式形式，却不能直接说明税收、司法或地方认同。对同一条约，应分别追问文本传本、使节往返、边市、堡寨和后续摩擦；对同一征服，应分别追问都城、乡村、边区和海域的时间差。这样写，才能避免把多政权世界误读为单一中枢的扩张史。

南下、崖山与海上远征 的意义正在于使、信息和人群流动短暂可见。材料不足之处，本书不补造完整场景，而以来源差异保留疑问；这既是对战争受难者和地方行动者的谨慎，也是对账本能够承担范围的界定。

\subsection{河患、红巾与北撤}

河决、起事、北伐和大都失守。战争与外交在，首先改变的往往不是地图颜色，而是道路能否通行、军粮由谁转运、城镇如何开闭、居民往何处迁徙以及地方首领同哪一方交涉。灾害不能单因解释覆亡 因此，征服和边界不应被写成一次战役之后立即完成的事实，而应被追踪为军事行动、环境约束和行政接收不断相互调整的过程。

这里的基本证据来自《元史》、\textit{《元朝秘史》}、波斯文和高丽材料、条格、碑刻与环境考古。不同政权的本纪、编年和战报常同意战事发生，却在军数、斩获、攻守责任和谈判条件上彼此冲突。军队报告天然服务奖惩和合法性，后出的亡国或胜利叙事又容易把复杂选择压缩成必然结果。故本书只把能够交叉确认的时间、地点和大致后果写作稳定节点；具体兵力、伤亡、城内情形和个人动机，均保留为需要多方材料核对的问题。

环境并非战争故事的背景布景。河流、山地、草场、渡口、运河、疫病、旱涝与季节决定补给速度，也影响逃亡、安置和复耕的可能。城址、古河道、道路遗存和区域考古能约束文本中的路线，却不能单独量化一次行军或一场灾害的损失。把地理条件同财政和文书联系起来，才可解释同一军令何以在不同地区产生不同的执行结果。

迁徙者、降附者、俘虏、军户、商人和寺院人员并不是被动的战争数字。他们在文书、墓志、碑刻、地方志或宗教材料中偶尔留下姓名与去向，但保存极不均衡。一个家族、一座城市或一批俘虏不能代表全体人口；同样，史书的沉默也不能被理解为没有流动。叙事应分开国家宣布的接收、军队实际控制和居民长期生活的改变。

外交语言也有自己的限制。册封、贡赐、兄弟之国、盟约、降附等称谓说明谈判的仪式形式，却不能直接说明税收、司法或地方认同。对同一条约，应分别追问文本传本、使节往返、边市、堡寨和后续摩擦；对同一征服，应分别追问都城、乡村、边区和海域的时间差。这样写，才能避免把多政权世界误读为单一中枢的扩张史。

河决、起事、北伐和大都失守 的意义正在于使、信息和人群流动短暂可见。材料不足之处，本书不补造完整场景，而以来源差异保留疑问；这既是对战争受难者和地方行动者的谨慎，也是对账本能够承担范围的界定。

\subsection{路线、气候与迁徙的证据限度}

战争中的迁徙并不总以官方名册出现。军队调动、城池易手、疫病、河流改道和粮道断绝，都可能迫使家庭暂时离开、改变职业或进入新的依附关系；但留下记录的往往只是使节、将领、官员或能够立碑的家族。对流民、俘虏和降附者的数字，应先辨它是战报、赈济统计还是后来的史家概括，不能把不同场景的数目相加。

同样，环境材料并不为文本提供一幅精确的背景地图。堤防遗址、沉积层、城址废弃和道路分期可以检验某些地理约束，却需要将始建、维修、毁弃和再利用分开。气候、灾害与战争之间可能互相放大，但相关并非单一因果。本书以路线、补给和安置的具体问题连接它们，而不把自然条件写成决定政权兴亡的宿命。

史书的沉默不应被填作无事发生：在没有地方证据的地段，只能确认军政记录所见的动作，不能断言居民是否顺从、逃离或获益。


